\subsection{Inheritance Without a Superclass Declaration}
\label{sec:original-inherit}

The control-form projects changed which expressions run. The original
\filepath{examples/misc/inherit.abclr2} now uses the same reflective
access to choose behavior from another object.
Familiar OOP gives us the question: can an emu reuse a bird's
behavior while overriding its flight behavior?
ABCL/R2's answer here is an explicit \emph{script-acquisition protocol},
not a class declaration or a built-in superclass search.

\subsubsection{Start with the Application, Not the Generator}

Load the original file in a fresh session:
\begin{lstlisting}
(inc "examples/misc/inherit" :use-source t)
\end{lstlisting}
The file defines a macro, a generator, and a group. Loading it does not
create bird or emu; enter their definitions from its closing comment:
\begin{codeblock}
\begin{lstlisting}
[object bird
 (script
  (=> :has-feather? !t)
  (=> :can-fly? !t))
 (group inherit-group)]

[object emu
 (script
  (=> :can-fly? !nil)
  (=> Any @ Reply from Sender
      (inherit Any Reply Sender bird)))
 (group inherit-group)]
\end{lstlisting}
\end{codeblock}

Both objects belong to \code{inherit-group}, which arranges for the
custom metaobject generator to be used.
The bird answers two questions locally. The emu answers the flight
question locally and forwards every other request through \code{inherit}.
The source's sample results are:
\begin{lstlisting}
[bird <== :can-fly?]           ; => T
[emu <== :can-fly?]            ; => NIL
[emu <== :has-feather?]        ; => T
\end{lstlisting}
\code{nil} in the second result is a reply containing false, not a
missing reply. The local emu handler sends it explicitly.

Check the identities separately from those application results:
\begin{lstlisting}
(eq [den [meta emu]] emu)             ; => T
(eq [meta emu] emu)                   ; => NIL
\end{lstlisting}
Emu's metaobject is another object, and its denotation leads back to emu.
Neither operation is a superclass lookup. These identity checks and
the three application replies pass on the current image.
Here \code{[group-of emu]} returns \code{inherit-group}. The generator
exposes its creation argument as \code{Group} for identity creation,
retaining \code{Group-manager} for execution. MCP, Meta-Meta, Fault-Tolerant,
and the Time-Warp application generator use the same convention.
This is application ownership, not the group of its metaobject or inherited method source.

For non-reifying metaobjects, the generator assigns its known creation
group to the denoted identity before initialization. The host expression
\code{(setf (r2sys::object-data-gid [den Metaobject]) GMgr)} writes that
identity's group slot; it does not inspect or expose the NRM state.

For the feather request, \code{Any} becomes \code{:has-feather?}.
\code{Reply} holds the caller's reply destination and \code{Sender}
holds the request's sender. These are context fields, not extra elements
inside the keyword message.
The fallback asks for bird's matching script with that context.
It does not first send the ordinary question \code{:has-feather?}
to bird and then copy bird's answer.

\subsubsection{The Reflective Boundary}

The macro from the Lisp reading guide expands the fallback call into:
\begin{lstlisting}
(reify-args-to-meta :inherit Any Reply Sender bird)
\end{lstlisting}
\code{reify-args-to-meta} sends compiled representations of these four
argument expressions, with their evaluation context, to emu's metaobject.
That is the object controlling emu's script execution.
The keyword \code{:inherit} selects a handler added by the custom generator.
The representations are not yet the four application values.

Read the core original handler now:
\wholelisting[linerange={33-35,38-48}]{examples/misc/inherit.abclr2}

The handler proceeds in three stages:
\begin{enumerate}
  \item \code{:do-evlis} asks the evaluator to evaluate the argument
        representations in their original environment. Its continuation
        receives the list \code{[message reply sender inherit-source]}.
  \item It asks the source's metaobject to find a script matching that
        message, reply destination, and sender. A successful reply contains
        pattern bindings and an executable script representation.
  \item It extends the current environment with those bindings and asks
        the evaluator to execute the acquired script in the current
        execution context, forwarding the meta-level continuation \code{C}.
\end{enumerate}
\code{Env} supplies the current bindings; \code{Id} identifies the
execution context used for \code{Me}; \code{Gid} supplies its group context;
\code{Eval} names the evaluator.
The acquired feather script sends \code{t} to the application's
\code{Reply}, answering the feather question. Separately, \code{C}
tells the reflective computation that the acquired expression has finished.
If no script is found, the failure branch warns and completes through
\code{C} with \code{nil}. It does not thereby send an application failure
reply to \code{Reply}.

\subsubsection{Reconstruct the Metaobject's Ordinary Dispatch}

Return to the generator's state and its normal-message handlers:
\wholelisting[firstline=23,lastline=30]{examples/misc/inherit.abclr2}

The application's initializers are described by \code{StateSpec}.
\code{LexEnv} names its outer lexical environment and \code{Scripts}
its script list. The group's execution components are \code{Eval}
and \code{GMgr}. The dotted \code{Extra} collects additional parameters.
The generated metaobject has its \emph{own} state: application queue,
application environment, scripts, evaluator, group manager, and mode.
The variable \code{state} is initially \code{nil}; it is filled by the
initialization code, not assumed to be the application's data already.

Read lines 49--71 of the original file with this state machine in mind:
\begin{enumerate}[nosep]
  \item \code{[:message Message Reply Sender]} queues one request.
        A dormant metaobject marks itself active and schedules \code{:begin}.
  \item \code{:begin} dequeues the request and finds its first matching
        script. The evaluator executes it in an environment containing
        the pattern bindings over the application's state.
  \item The execution continuation sends \code{:end} to the metaobject.
        That handler starts another queued request, or becomes dormant.
\end{enumerate}
\code{Me} in these metaobject handlers denotes the metaobject itself.
\code{[den Me]} obtains the application object it describes.
The custom \code{:inherit} handler is added to this machinery; it does
not replace the ordinary application inbox by a superclass hierarchy.

\subsubsection{Initialization and Reflective Queries}

Lines 73--92 expose lookup and update operations such as \code{:state},
\code{:scripts}, \code{:find-script}, and \code{:set-evaluator}.
In particular, the three-field \code{:find-script} query preserves sender
and reply information; the one-field query uses \code{nil} for those fields.
That difference matters for patterns that depend on the sender.

Finally, lines 95--104 evaluate state initializers with \code{:do-evlis}.
\code{stateSpec-varNames} supplies names and \code{pairlis} associates
them with the resulting values. \code{env-gen} creates an environment
with those bindings and the lexical parent \code{LexEnv}.
The generator sets that environment on the new metaobject and replies
with the metaobject to its creation caller.
Read initialization last, once you know what its state must support.

\paragraph{Understanding check, with answers.}
Why does emu's \code{:can-fly?} not consult bird?
Its specific local pattern is first and matches, so the catch-all never runs.
Would moving the catch-all first preserve that override? No.
Does the example copy bird's state or permanently attach bird's scripts
to emu? No: it finds a script for each inherited request and executes it
using the current environment and context. It is not a guarantee of the
field or superclass semantics of another OOP language.
Why is \code{from Sender} retained? Acquired script matching may depend
on who sent the original request, not just its payload.

\paragraph{Reading milestone.}
You should now be able to locate the application reply, the expression
completion, the queue transition, and the creation reply independently
in this file. The original baseline checks the three bird/emu queries
without changing any original expression.
That success does not certify arbitrary scripts, runtime script
mutation, or every use of the generator.
